\chap{自适应验证：语义规范、代码测试与真实使用}
\label{sec:validation}
\subsection{验证要支持什么判断}
验证的价值在于减少对当前软件行为的不确定性。一次检查应能够说明它针对什么任务或契约、使用了什么内容与环境、观察到什么结果。对于持续变化的系统，验证结果具有范围和条件，不能把过去某次成功无条件延续为当前整个系统的可靠性。

代码式测试适合反复检查稳定且可精确表达的规则，例如请求幂等性、状态转移约束和权限边界。语义规范说明为什么需要这些规则，以及用户任务怎样才算成功。TAU 则让 Agent 从真实入口完成任务，检查多个模块组合之后的实际行为，并发现难以提前写成断言的体验问题。三者共同支持判断，维护方式也需要随变化调整。

真实任务基准为这一方向提供了参考。SWE-bench 将仓库与真实问题联系起来，要求系统修改代码解决问题\cite{swebench}；WebArena 与 OSWorld 分别构造网页和真实计算机环境中的任务评价\cite{webarena,osworld}。这些工作表明，评价可以围绕环境中的任务结果组织。本文把这种任务视角用于软件开发中的持续验证，并关注验证发现如何进入修复和交付。

\subsection{以完整用户任务组织 TAU}
TAU 的选择单位可以是用户旅程。例如，创建一个长时间任务、观察进度、提出取消、确认资源释放并查看已保留结果，形成一条完整任务。它自然经过多个模块，也能暴露单元测试难以发现的状态解释问题。任务无需与每个设计点一一对应。

运行前需要明确初始条件、可使用入口、成功含义和可观察结果。若测试目标是普通用户体验，Agent 就不应依赖隐藏数据库修改来完成任务；必要的后台核对可以作为独立证据，同时说明它与用户可见体验的区别。测试数据应使关键分支可重复，例如让执行器延迟停止，观察界面是否过早显示终态。

Agent 的适应能力可能掩盖产品问题。它可能通过猜测按钮含义、刷新页面或绕过不清楚的入口最终完成任务。因此，TAU 记录应同时包含结果和过程中的摩擦：反复尝试、等待缺乏反馈、术语不一致、错误提示难以理解等。完成任务不能独自证明普通用户能够顺畅使用。

\begin{table}[htbp]\centering\small
\caption{取消任务旅程中的不同验证责任}\label{tab:tau}
\begin{tabularx}{\linewidth}{p{2.5cm}YY}\toprule
行为 & 低成本代码检查 & 真实使用需要观察的内容\\\midrule
重复点击取消 & 幂等键与状态不重复推进 & 用户是否理解请求已被接收\\
执行器延迟停止 & 非终态不被提前写为终态 & 等待反馈是否明确、是否误导\\
取消与完成并发 & 实际终态与版本一致 & 页面如何解释任务已经完成\\
保留中间结果 & 输出索引与访问规则 & 用户能否找到并理解保留内容\\\bottomrule
\end{tabularx}\end{table}

\subsection{证据在变化后的有效范围}
可将一条验证证据表示为
\begin{equation}
 e=(s,v,c,o,a),
\end{equation}
其中 $s$ 为检查范围，$v$ 为实际内容或版本，$c$ 为环境条件，$o$ 为观察结果，$a$ 为结论依赖的假设。这个表示用于解释需要记录哪些事实，具体系统不必采用相同的数据结构。一次变化 $\Delta$ 是否使证据失效，要看它是否影响范围或假设，而不只看时间先后。

例如，取消提示文案变化通常不会使执行器资源释放测试失效；执行器停止逻辑变化则可能影响该结论，即使前端代码完全没变。依赖记录可以辅助选择，但不必假定它完备。当影响范围无法可靠限定时，扩大检查是合理选择。维护证据范围本身也有成本，需要与直接运行测试的成本比较。

\fig{07-evidence}{局部证据复用示例。变化只使相关检查需要更新，其他范围的有效结论仍可保留；跨模块组合行为可以通过选定的真实任务补查。}{evidence}

在多人或多 Agent 持续开发中，测试执行期间仍可能出现新变化。结果可以说明它实际检查的内容，再对后续相关差异补查。让任何新修改都导致整轮测试作废，会形成难以结束的追赶。相反，把实际内容、依赖和未检查范围记录清楚，有助于逐步建立判断。

\subsection{何时扩大验证}
验证策略可以考虑变化重要性、已有覆盖、近期故障、组合未知和执行成本。单一局部变化可以只检查相关责任范围；多个模块连续改变同一用户任务时，值得检查组合路径；基础协议变化且影响难以限定时，可以扩大到更广范围。上述情境是判断线索，不构成固定触发阈值。

全量代码测试和 TAU 的范围可以分别选择。某项目全量代码测试运行很快，精细分析用例反而更贵；而交互任务可能需要大量模型调用，适合优先选择核心旅程和变化最显著的路径。系统应比较实际成本，而不机械地把“更精准选择”视为更高效。

若以 $V(q)$ 表示检查 $q$ 的预期判断价值、$C(q)$ 表示其资源成本，可以把优先级理解为在给定预算下选择有价值的检查组合。由于价值估计具有不确定性，实际系统应保留选择依据并根据结果更新经验。形式化表达帮助明确问题，不能替代工程证据。

\fig{08-validation}{自适应验证的持续过程。范围判断、检查执行、结果解释、修复与证据更新相互连接；失败后的处理继续由相关责任成员完成。}{validation}

\subsection{三种不同变化下的投入选择}
假设第一轮变化仅调整任务列表中的提示文案，已有状态契约和执行逻辑均未改变。可以核对文字含义、布局和状态映射，并复用执行器相关证据。若提示文案改变了用户对终态的理解，则需要增加一项聚焦的体验检查；“只改文字”本身并不能保证没有行为影响。

第二轮变化同时调整任务服务的重试策略和执行器响应方式。各模块局部检查通过后，仍存在组合条件未知，例如重试是否造成重复终态事件。这时以取消旅程组织一次相关 TAU，并核对事件幂等，可以更直接地获得判断依据。

第三轮变化替换了多个基础组件，依赖映射也不够完整。精细选测的依据变弱，可以运行较广的代码回归，同时选择几条核心用户任务。若全量代码检查成本很低，直接执行可能更合适。三种情境说明策略取决于需要消除的未知，而不是固定发布次数或固定日程。

\subsection{从失败到修复完成}
检查失败后，首先要区分产品缺陷、过时断言、环境问题以及测试操作本身的错误。区分依据可以包括复现、独立观察和当前规范。执行者不应为了得到通过结果任意放宽测试；也不应把所有失败都当成产品问题，造成无关修改。

明确原因之后，相关成员完成修复，重跑失败项和受影响范围，再通过既定交付路径送达。若问题反复出现且已经能够精确描述，可以增加低成本代码检查，减少以后重复支付探索成本。仍依赖界面感受和任务理解的问题，继续通过适合的真实使用方法核对。

修复跟踪需要停止条件。当前问题已得到解释、修复已核对、剩余未知被明确记录时，可以结束本轮工作。新出现的独立问题可以形成后续任务。无限扩展检查范围会消耗预算并延迟成果送达；过早结束则会让“发现问题”与“完成修复”脱节。

\subsection{让验证任务与结果可以交接}
一次 TAU 可以使用简短任务说明，把目标、初始条件、允许入口、成功含义和证据要求交给执行者。任务说明应描述用户需要完成的事情，避免预先写死每一步点击，从而允许发现新的有效路径和实际摩擦。对于必须核对的边界条件，则应明确给出相应情境。

\begin{table}[htbp]\centering\small
\caption{可交接的验证任务与结果记录示例}\label{tab:tau-record}
\begin{tabularx}{\linewidth}{p{2.4cm}Y}\toprule
字段 & 取消任务示例\\\midrule
目标与情境 & 创建长时间任务，在执行器延迟停止时取消，理解实际终态并查看保留结果\\
允许入口 & 用户界面完成操作；后台观察单独记录，不用于替代用户步骤\\
接受条件 & 请求已接收可见，停止前不显示停止完成，终态及结果可访问\\
观察材料 & 实际内容与环境、关键状态、操作序列、相关运行证据\\
结果解释 & 哪些条件成立，哪里需要猜测或绕路，失败涉及哪个对象\\
后续处理 & 产品、测试或环境问题的判断；修复责任与必要复查范围\\\bottomrule
\end{tabularx}\end{table}

结果记录可以同时服务于开发和汇报。模块成员需要复现条件与证据，人类读者可能只需要理解问题如何影响任务。二者共享同一事实入口，分别选择适合的表达。若结果仅保存为一段无法定位对象的总结，后续修复会再次支付背景重建成本。

任务和结果的交接也要保留未完成项。执行环境中断时，记录已经完成的步骤与当前状态，恢复后从可信位置继续；若初始条件已经变化，则重新判断哪些步骤需要重做。这样可以减少简单重跑，同时避免沿用已经不成立的结论。

\subsection{验证的组织与可靠性}
验证职责可以由 Leader 委派给专门成员，也可以由模块成员与独立体验者合作承担。重要的是检查方式能接触不同证据，并能把问题送回合适的责任范围。固定的角色名称无法保证独立性，多个成员共享同一个错误假设时仍可能共同遗漏问题。

评价验证策略时，应观察遗漏的重要缺陷、误报带来的返工、修复时间和每次检查获得的信息。只报告测试数量或通过率可能产生误导：大量简单检查可以全部通过，而核心用户任务仍然不可用。对未覆盖范围保持清楚认识，能帮助人和团队作出更合理的后续选择。

TAU 也需要接受人类使用的校正。真实用户可能采用 Agent 没有想到的路径，也可能无法理解 Agent 轻易推断出的术语。可以通过小规模用户任务观察检查这种差异，再调整验证任务与界面表达。由此，验证过程本身成为持续学习的对象。
